%% Преамбула TeX-файла

% 1. Стиль и язык
\documentclass[utf8x, 12pt]{G7-32} % Стиль (по умолчанию будет 14pt)

% Остальные стандартные настройки убраны в preamble.inc.tex.
\sloppy

% Настройки стиля ГОСТ 7-32
% Для начала определяем, хотим мы или нет, чтобы рисунки и таблицы нумеровались в пределах раздела, или нам нужна сквозная нумерация.
\EqInChapter % формулы будут нумероваться в пределах раздела
\TableInChapter % таблицы будут нумероваться в пределах раздела
\PicInChapter % рисунки будут нумероваться в пределах раздела
\usepackage{slashbox}

% Добавляем гипертекстовое оглавление в PDF
\usepackage[
bookmarks=true, colorlinks=true, unicode=true,
urlcolor=black,linkcolor=black, anchorcolor=black,
citecolor=black, menucolor=black, filecolor=black,
]{hyperref}

% Изменение начертания шрифта --- после чего выглядит таймсоподобно.
% apt-get install scalable-cyrfonts-tex

\IfFileExists{cyrtimes.sty}
    {
        \usepackage{cyrtimespatched}
    }
    {
        % А если Times нету, то будет CM...
    }

\usepackage{graphicx}   % Пакет для включения рисунков

% С такими оно полями оно работает по-умолчанию:
% \RequirePackage[left=20mm,right=10mm,top=20mm,bottom=20mm,headsep=0pt]{geometry}
% Если вас тошнит от поля в 10мм --- увеличивайте до 20-ти, ну и про переплёт не забывайте:
\geometry{right=20mm}
\geometry{left=30mm}


% Пакет Tikz
\usepackage{tikz}
\usetikzlibrary{arrows,positioning,shadows}

% Произвольная нумерация списков.
\usepackage{enumerate}

% ячейки в несколько строчек
\usepackage{multirow}

% itemize внутри tabular
\usepackage{paralist,array}

% Центрирование подписей к плавающим окружениям
\usepackage[justification=centering]{caption}


% Настройки листингов.
\ifPDFTeX
% 8 Листинги

\usepackage{listings}
\usepackage{wrapfig}
% Значения по умолчанию
\lstset{
  basicstyle= \footnotesize,
  breakatwhitespace=true,% разрыв строк только на whitespacce
  breaklines=true,       % переносить длинные строки
%   captionpos=b,          % подписи снизу -- вроде не надо
  inputencoding=koi8-r,
  numbers=left,          % нумерация слева
  numberstyle=\footnotesize,
  showspaces=false,      % показывать пробелы подчеркиваниями -- идиотизм 70-х годов
  showstringspaces=false,
  showtabs=false,        % и табы тоже
  stepnumber=1,
  tabsize=4,              % кому нужны табы по 8 символов?
  frame=single,
  escapeinside={(*}{*)}, %выделение
  literate={а}{{\selectfont\char224}}1
  {б}{{\selectfont\char225}}1
  {в}{{\selectfont\char226}}1
  {г}{{\selectfont\char227}}1
  {д}{{\selectfont\char228}}1
  {е}{{\selectfont\char229}}1
  {ё}{{\"e}}1
  {ж}{{\selectfont\char230}}1
  {з}{{\selectfont\char231}}1
  {и}{{\selectfont\char232}}1
  {й}{{\selectfont\char233}}1
  {к}{{\selectfont\char234}}1
  {л}{{\selectfont\char235}}1
  {м}{{\selectfont\char236}}1
  {н}{{\selectfont\char237}}1
  {о}{{\selectfont\char238}}1
  {п}{{\selectfont\char239}}1
  {р}{{\selectfont\char240}}1
  {с}{{\selectfont\char241}}1
  {т}{{\selectfont\char242}}1
  {у}{{\selectfont\char243}}1
  {ф}{{\selectfont\char244}}1
  {х}{{\selectfont\char245}}1
  {ц}{{\selectfont\char246}}1
  {ч}{{\selectfont\char247}}1
  {ш}{{\selectfont\char248}}1
  {щ}{{\selectfont\char249}}1
  {ъ}{{\selectfont\char250}}1
  {ы}{{\selectfont\char251}}1
  {ь}{{\selectfont\char252}}1
  {э}{{\selectfont\char253}}1
  {ю}{{\selectfont\char254}}1
  {я}{{\selectfont\char255}}1
  {А}{{\selectfont\char192}}1
  {Б}{{\selectfont\char193}}1
  {В}{{\selectfont\char194}}1
  {Г}{{\selectfont\char195}}1
  {Д}{{\selectfont\char196}}1
  {Е}{{\selectfont\char197}}1
  {Ё}{{\"E}}1
  {Ж}{{\selectfont\char198}}1
  {З}{{\selectfont\char199}}1
  {И}{{\selectfont\char200}}1
  {Й}{{\selectfont\char201}}1
  {К}{{\selectfont\char202}}1
  {Л}{{\selectfont\char203}}1
  {М}{{\selectfont\char204}}1
  {Н}{{\selectfont\char205}}1
  {О}{{\selectfont\char206}}1
  {П}{{\selectfont\char207}}1
  {Р}{{\selectfont\char208}}1
  {С}{{\selectfont\char209}}1
  {Т}{{\selectfont\char210}}1
  {У}{{\selectfont\char211}}1
  {Ф}{{\selectfont\char212}}1
  {Х}{{\selectfont\char213}}1
  {Ц}{{\selectfont\char214}}1
  {Ч}{{\selectfont\char215}}1
  {Ш}{{\selectfont\char216}}1
  {Щ}{{\selectfont\char217}}1
  {Ъ}{{\selectfont\char218}}1
  {Ы}{{\selectfont\char219}}1
  {Ь}{{\selectfont\char220}}1
  {Э}{{\selectfont\char221}}1
  {Ю}{{\selectfont\char222}}1
  {Я}{{\selectfont\char223}}1
}

% Стиль для псевдокода: строчки обычно короткие, поэтому размер шрифта побольше
\lstdefinestyle{pseudocode}{
  basicstyle=\small,
  keywordstyle=\color{black}\bfseries\underbar,
  language=Pseudocode,
  numberstyle=\footnotesize,
  commentstyle=\footnotesize\it
}

% Стиль для обычного кода: маленький шрифт
\lstdefinestyle{realcode}{
  basicstyle=\scriptsize,
  numberstyle=\footnotesize
}

% Стиль для коротких кусков обычного кода: средний шрифт
\lstdefinestyle{simplecode}{
  basicstyle=\footnotesize,
  numberstyle=\footnotesize
}

% Стиль для BNF
\lstdefinestyle{grammar}{
  basicstyle=\footnotesize,
  numberstyle=\footnotesize,
  stringstyle=\bfseries\ttfamily,
  language=BNF
}

% Определим свой язык для написания псевдокодов на основе Python
\lstdefinelanguage[]{Pseudocode}[]{Python}{
  morekeywords={each,empty,wait,do},% ключевые слова добавлять сюда
  morecomment=[s]{\{}{\}},% комменты {а-ля Pascal} смотрятся нагляднее
  literate=% а сюда добавлять операторы, которые хотите отображать как мат. символы
    {->}{\ensuremath{$\rightarrow$}~}2%
    {<-}{\ensuremath{$\leftarrow$}~}2%
    {:=}{\ensuremath{$\leftarrow$}~}2%
    {<--}{\ensuremath{$\Longleftarrow$}~}2%
}[keywords,comments]

% Свой язык для задания грамматик в BNF
\lstdefinelanguage[]{BNF}[]{}{
  morekeywords={},
  morecomment=[s]{@}{@},
  morestring=[b]",%
  literate=%
    {->}{\ensuremath{$\rightarrow$}~}2%
    {*}{\ensuremath{$^*$}~}2%
    {+}{\ensuremath{$^+$}~}2%
    {|}{\ensuremath{$|$}~}2%
}[keywords,comments,strings]

% Подписи к листингам на русском языке.
\renewcommand\lstlistingname{\cyr\CYRL\cyri\cyrs\cyrt\cyri\cyrn\cyrg}
\renewcommand\lstlistlistingname{\cyr\CYRL\cyri\cyrs\cyrt\cyri\cyrn\cyrg\cyri}

\else
\usepackage{local-minted}
\fi

% Полезные макросы листингов.
% Любимые команды
\newcommand{\Code}[1]{\textbf{#1}}


\begin{document}

\frontmatter % выключает нумерацию ВСЕГО; здесь начинаются ненумерованные главы: реферат, введение, глоссарий, сокращения и прочее.

% Команды \breakingbeforechapters и \nonbreakingbeforechapters
% управляют разрывом страницы перед главами.
% По-умолчанию страница разрывается.

% \nobreakingbeforechapters
% \breakingbeforechapters

% Также можно использовать \Referat, как в оригинале
%\begin{abstract}
%	Титульный лист. Эта страница нужна мне, чтобы не сбивалась нумерация страниц
%	\cite{Dh}
%	\cite{Bayer}
%	\cite{Habr1}
%	\cite{Noise_func}
%	\cite{Ulich}

%Это пример каркаса расчётно-пояснительной записки, желательный к использованию в РПЗ проекта по курсу РСОИ.

%Данный опус, как и более новые версии этого документа, можно взять по адресу (\url{https://github.com/rominf/latex-g7-32}).

%Текст в документе носит совершенно абстрактный характер.
%\end{abstract}
% НАЧАЛО ТИТУЛЬНОГО ЛИСТА
\begin{center}
	\hfill \break
	\textit{
		\normalsize{Государственное образовательное учреждение высшего профессионального образования}}\\ 
	
	\textit{
		\normalsize  {\bf  «Московский государственный технический университет}\\ 
		\normalsize  {\bf имени Н. Э. Баумана»}\\
		\normalsize  {\bf (МГТУ им. Н.Э. Баумана)}\\
	}
	\noindent\rule{\textwidth}{2pt}
	\hfill \break
	\noindent
	\makebox[0pt][l]{ФАКУЛЬТЕТ}%
	\makebox[\textwidth][c]{«Информатика и системы управления»}%
	\\
	\noindent
	\makebox[0pt][l]{КАФЕДРА}%
	\makebox[\textwidth][r]{«Программное обеспечение ЭВМ и информационные технологии»}%
	\\
	\hfill\break
	\hfill \break
	\hfill \break
	\hfill \break
	\normalsize{\bf Р А С Ч Ё Т Н О - П О Я С Н И Т Е Л Ь Н А Я\space\space З А П И С К А}\\
	\normalsize{\bf к курсовой работе на тему:}\\
	\hfill \break
	\large{Метод автоматизированного поиска и вывода изоморфных контекстно-зависимых визуальных отображений на изображении}\\
	\hfill \break
	\hfill \break
	\hfill \break
	\hfill \break
	\hfill \break	
	\normalsize {
		\noindent
		\makebox[0pt][l]{Студент}%
		\makebox[\textwidth][c]{Хмельщиков А.Ю.}%
		\makebox[0pt][r]{{$\underset{\text{(Подипсь, дата)}}{\underline{\hspace{8cm}}}$ \space И.О. Фамилия}}
	}\\
	\hfill \break	
	\normalsize {
		\noindent
		\makebox[0pt][l]{Руководитель курсового проекта}%
		\makebox[\textwidth][c]{ ~~~~~~~~      Оленев А.А.}%
		\makebox[0pt][r]{{$\underset{\text{(Подпись, дата)}}{\underline{\hspace{6cm}}}$ \space И.О. Фамилия}}
	}
	\hfill \break
	\hfill \break
	\hfill \break
	\hfill \break
\end{center}
\hfill \break
\hfill \break
\begin{center} Москва 2017 \end{center}

\thispagestyle{empty} % 
% КОНЕЦ ТИТУЛЬНОГО ЛИСТА


%%% Local Variables: 
%%% mode: latex
%%% TeX-master: "rpz"
%%% End: 

\tableofcontents

%
%%% Local Variables:
%%% mode: latex
%%% TeX-master: "rpz"
%%% End:

%%
%%% Local Variables:
%%% mode: latex
%%% TeX-master: "rpz"
%%% End:


Intro: The quick brown fox jumps over the lazy dog.  The quick brown fox jumps over the lazy dog.  The quick brown fox jumps over the lazy dog.  The quick brown fox jumps over the lazy dog.  The quick brown fox jumps over the lazy dog.The quick brown fox jumps over the lazy dog.  The quick brown fox jumps over the lazy dog.  The quick brown fox jumps over the lazy dog.  The quick brown fox jumps over the lazy dog.  The quick brown fox jumps over the lazy dog.The quick brown fox jumps over the lazy dog.  The quick brown fox jumps over the lazy dog.  The quick brown fox jumps over the lazy dog.  The quick brown fox jumps over the lazy dog.  The quick brown fox jumps over the lazy dog.The quick brown fox jumps over the lazy dog.  The quick brown fox jumps over the lazy dog.  The quick brown fox jumps over the lazy dog.  The quick brown fox jumps over the lazy dog.  The quick brown fox jumps over the lazy dog.The quick brown fox jumps over the lazy dog.  The quick brown fox jumps over the lazy dog.  The quick brown fox jumps over the lazy dog.  The quick brown fox jumps over the lazy dog.  The quick brown fox jumps over the lazy dog.The quick brown fox jumps over the lazy dog.  The quick brown fox jumps over the lazy dog.  The quick brown fox jumps over the lazy dog.  The quick brown fox jumps over the lazy dog.  The quick brown fox jumps over the lazy dog.The quick brown fox jumps over the lazy dog.  The quick brown fox jumps over the lazy dog.  The quick brown fox jumps over the lazy dog.  The quick brown fox jumps over the lazy dog.  The quick brown fox jumps over the lazy dog.The quick brown fox jumps over the lazy dog.  The quick brown fox jumps over the lazy dog.  The quick brown fox jumps over the lazy dog. The quick brown fox jumps over the lazy dog.  The quick brown fox jumps over the lazy dog.

\mainmatter % это включает нумерацию глав и секций в документе ниже

\chapter{Аналитический раздел}
\label{cha:analysis}
%
% % В начале раздела  можно напомнить его цель
%
\section{Предыдущие работы}
\subsection{Suggestive Sketching}
Drawing is a common and yet challenging human activity. Significant research efforts have been devoted to design suggestive or guided drawing systems that use various forms of data to help users. For example, portrait sketching can be assisted by crowd- sourced sketches \cite{bib1} or analyzed face data \cite{bib1}. To help users draw a larger collection of objects, Lee et al. \cite{bib1} interactively guides users’ progress by displaying shadows extracted from web images, while Iarussi et al. \cite{bib1} provide structural guidances based on artistic and geometric principles. Besides static images or drawings, it is also possible to guide novices through recorded workflows of experienced painters \cite{bib1}.
Our method follows this line of work, but uses past drawing workflows from the same user to guide future drawings.

\subsection{Data-driven Sketching}
Although the guided systems can help users draw easier, they still have to deal with manual repetitions. To help ameliorate such tedious process, many systems have been designed to automate the authoring of repetitive drawings via data driven computation. One prime example is the detailed textures or patterns  \cite{bib1} \cite{bib1} \cite{bib1} \cite{bib1} whose creation can fit well with the traditional copy-and-paste interaction styles. Kazi et al.  \cite{bib1} \cite{bib1} provide friendly gesture controls to help create both static and dynamic elements.
Our method also follows a data-driven approach to help creating repetitions, but uses dynamic workflows rather than static patterns or animated sprites.

\subsection{Using Workflows}
Chimera \cite{bib1} is an old system for recording and editing graphical histories, specifically the repetitive operations. Recent years have seen the rise of methods that utilize workflows in various forms, such as exploration \cite{bib1}, visualization  \cite{bib1} \cite{bib1} \cite{bib1} \cite{bib1}, stylization \cite{bib1}, beautification \cite{bib1}, tutorial  \cite{bib1} \cite{bib1}, and revision \cite{bib1}. These methods often rely on pre-recorded workflows, which, when not available, may also be recovered through a certain extent of analysis  \cite{bib1} \cite{bib1} \cite{bib1} \cite{bib1}. Recently, Nancel et al. \cite{bib1} provides a comprehensive survey of different conceptual models for workflow analysis.
We follow this line of work, but focus on the analysis and synthesis of dynamic workflows for suggestive painting repetitions.


% Обратите внимание, что включается не ../dia/..., а inc/dia/...
% В Makefile есть соответствующее правило для inc/dia/*.pdf, которое
% берет исходные файлы из ../dia в этом случае.
%$  $

%%% Local Variables:
%%% mode: latex
%%% TeX-master: "rpz"
%%% End:

\chapter{Конструкторский раздел}
\label{cha:design}
\section{Общая архитектура приложения}
\begin{figure}[h!]
	\centering
	\includegraphics[width=0.6\textwidth]{img/fox.jpg}
	\caption{Общая архитектура приложения}
	\label{fig:spire09}
\end{figure}
The quick brown fox jumps over the lazy dog.The quick brown fox jumps over the lazy dog.The quick brown fox jumps over the lazy dog.The quick brown fox jumps over the lazy dog.The quick brown fox jumps over the lazy dog.The quick brown fox jumps over the lazy dog.The quick brown fox jumps over the lazy dog.


The quick brown fox jumps over the lazy dog.The quick brown fox jumps over the lazy dog.The quick brown fox jumps over the lazy dog.The quick brown fox jumps over the lazy dog.The quick brown fox jumps over the lazy dog.The quick brown fox jumps over the lazy dog.The quick brown fox jumps over the lazy dog.The quick brown fox jumps over the lazy dog.The quick brown fox jumps over the lazy dog.The quick brown fox jumps over the lazy dog.The quick brown fox jumps over the lazy dog.

%%% Local Variables:

%%% mode: latex
%%% TeX-master: "rpz"
%%% End:
%--количество цветов
%||количество пикселей
\chapter{Технологический раздел}
\section{Выбор  языка программирования}
Для реализации  был выбран язык  С++ по следующими причинам:
\begin{enumerate}
	 \item The quick brown fox jumps over the lazy dog.
	 \item Компилируемый язык со статической типизацией. 
	 \item Сочетание высокоуровневых и низкоуровневых средств.
	 \item Реализация ООП.
	 \item Наличие удобной стандартной библиотеки шаблонов
	 \end{enumerate}
\section{Выбор вспомогательных библиотек}
Для реализации программы была выбрана библиотека Qt.
\begin{enumerate}
	\item The quick brown fox jumps over the lazy dog.
	\item The quick brown fox jumps over the lazy dog.
\end{enumerate}




%%% Local Variables:
%%% mode: latex
%%% TeX-master: "rpz"
%%% End:

%%\chapter{Исследовательский раздел}
%\chapter{Заключение}
%\label{cha:research}
%\section{Время дизеринга раличных алгоритмов}

%В данном разделе проводятся вычислительные эксперименты.

%%% Local Variables:
%%% mode: latex
%%% TeX-master: "rpz"
%%% End:



\backmatter %% Здесь заканчивается нумерованная часть документа и начинаются ссылки и
            %% заключение

\Conclusion % заключение к отчёту

\begin{figure}
The quick brown fox jumps over the lazy dog.The quick brown fox jumps over the lazy dog.The quick brown fox jumps over the lazy dog.The quick brown fox jumps over the lazy dog.The quick brown fox jumps over the lazy dog.The quick brown fox jumps over the lazy dog.The quick brown fox jumps over the lazy dog.The quick brown fox jumps over the lazy dog.The quick brown fox jumps over the lazy dog.The quick brown fox jumps over the lazy dog.The quick brown fox jumps over the lazy dog.The quick brown fox jumps over the lazy dog.
\end{figure}
%%% Local Variables: 
%%% mode: latex
%%% TeX-master: "rpz"
%%% End: 


% % Список литературы при помощи BibTeX
% Юзать так:
%
% pdflatex rpz
% bibtex rpz
% pdflatex rpz

\bibliographystyle{gost780u}
\bibliography{rpz}


%%% Local Variables: 
%%% mode: latex
%%% TeX-master: "rpz"
%%% End: 


%\appendix   % Тут идут приложения

%%\chapter{Картинки}
%\label{cha:appendix1}

%\begin{figure}
%\centering
%\caption{Картинка в приложении. Страшная и ужасная.}
%\end{figure}

%%% Local Variables: 
%%% mode: latex
%%% TeX-master: "rpz"
%%% End: 

%%\chapter{Еще картинки}
%\label{cha:appendix2}

%\begin{figure}
%\centering
%\caption{Еще одна картинка, ничем не лучше предыдущей. Но %надо же как-то заполнить место.}
%\end{figure}

%%% Local Variables: 
%%% mode: latex
%%% TeX-master: "rpz"
%%% End: 


\end{document}

%%% Local Variables:
%%% mode: latex
%%% TeX-master: t
%%% End:
